\newpage
\section{REMAINING TASKS}

Table~\ref{tab:impl_status} lists the work remaining after this progress report, consistent with the two-phase scope described throughout the Implementation chapter (Section~5): baseline and SOTA model reproduction on clean ASVspoof5 is complete, and the items below cover the ablation completion and RTC-aware pipeline integration planned for the final report.

\renewcommand{\arraystretch}{1.5}
\setlength{\tabcolsep}{8pt}
\begin{longtable}{>{\raggedright\arraybackslash}p{4.2cm} >{\raggedright\arraybackslash}p{8.5cm}}
\caption{Remaining Tasks and Execution Plan} \label{tab:impl_status} \\
\hline
\textbf{Remaining Task} & \textbf{Plan to Execute} \\
\hline
\endfirsthead
\hline
\textbf{Remaining Task} & \textbf{Plan to Execute} \\
\hline
\endhead
\hline
\endfoot
Investigate the RawNet2/AASIST evaluation-set EER gap and capture their missing metrics (Section~6.4) &
RawNet2's development-set EER (5.87\%) is now captured and is competitive with the classical baselines, yet its evaluation-set EER (36.04\%) remains far higher, pointing toward a generalization gap to the unseen attack types in the evaluation partition rather than under-training; its precision, recall, F1, and ROC-AUC are still missing. AASIST's evaluation-set EER (29.12\%) is known, but its development-set EER and full supporting metrics have not yet been captured at all. Confirm the generalization-gap hypothesis for RawNet2, diagnose AASIST from scratch, and capture all missing metrics for both. \\
RTC-aware distortion pipeline: codec round-trip, packet-loss simulation, post-processing (Section~4.8.1) &
Implement the codec round-trip (Opus, G.711, AAC via \texttt{ffmpeg}), the Gilbert--Elliott bursty packet-loss simulator, and the AGC/high-pass post-processing stage, gated by a clean-sample probability during training. \\
RTC transmission pipeline integration: sender, network, receiver subsystems (Sections~4.1--4.4) &
Implement the sender-side capture/encoding/packetization chain, the network impairment layer, and the receiver-side jitter-buffer/PLC/decoding/enhancement chain, reusing the codec and packet-loss models built for the augmentation pipeline above. \\
Detection service integration with RTC pipeline (Sections~4.5--4.6) &
Deploy the trained detection models behind the streaming API (gRPC/HTTP) described in Section~4.5, connect it to the receiver-side output, and validate chunked inference and temporal aggregation end-to-end. \\
RTC distortion robustness evaluation across codecs and packet-loss rates (Section~4.8) &
Re-evaluate all eight baseline models under the full codec $\times$ packet-loss-rate sweep, using the same EER-based protocol established in Sections~5.2.5 and~5.3.5, to quantify degradation relative to the clean-audio results in Section~6.4. \\
Real-time performance profiling: latency, throughput, resource utilization (Section~4.7) &
Measure inference latency, throughput, and memory/CPU/GPU utilization for each model behind the deployed detection service, using the metrics defined in Section~4.7.2. \\
\end{longtable}

The clean-audio baseline comparison established in this phase provides the reference operating points for each of the eight models (Section~6.4). These reference points will be used to measure the extent of RTC-induced performance degradation once the RTC-aware augmentation pipeline described in Section~4.8.1 is integrated, directly addressing the RTC distortion robustness research question defined in Section~4.8.
